\section{Bases de diseño}

\subsection{Elementos a fundar}

El diseño abarca los 30 elementos verticales que descargan al nivel de fundación: 21 muros de hormigón armado, resueltos mediante zapata corrida, y 9 columnas, resueltas mediante zapata aislada. Los muros se agrupan según su espesor y posición: los muros perimetrales presentan un espesor $e_w = 0,20$ m (designación MHA20), mientras que los muros interiores tienen $e_w = 0,25$ m (designación MHA25). Las columnas poseen una sección transversal de $30 \times 50$ cm. La totalidad de los elementos a fundar emplea hormigón G35. Dado el comportamiento esencialmente unidireccional del muro, su zapata corrida se verifica en una dirección, en tanto que las zapatas aisladas de columna se verifican de forma biaxial.

Los cinco muros perimetrales del sector, P1, PN, P13, Plb y P12, se resuelven mediante zapata excéntrica: al ubicarse en el límite del edificio, su base no puede desarrollarse hacia el exterior sin sobrepasar la línea de cierro, por lo que se diseñan con el muro en el borde de la zapata y la base proyectándose únicamente hacia el interior, conforme al criterio del curso para fundaciones de muros perimetrales. La metodología correspondiente se detalla en la sección~\ref{subsec:zapata_excentrica}.

El foso de ascensor, delimitado por los muros PH, P9 y PE, se ubica a una profundidad de 1,50 m bajo el nivel de fundación y requiere una losa de fondo, modelada como una placa apoyada en sus cuatro bordes según el criterio del curso.

\subsection{Materiales}
El hormigón de las fundaciones es G35, con resistencia a la compresión y módulo de elasticidad
\begin{equation}
f'_c = 350~\mathrm{kg/cm^2}, \qquad E_c = 15100\sqrt{350} = 282\,574~\mathrm{kg/cm^2},
\end{equation}
especificado según NCh170. El acero de refuerzo es A630-420H según NCh204, empleado en la totalidad de las fundaciones, con
\begin{equation}
f_y = 4200~\mathrm{kg/cm^2}, \qquad E_s = 2\,000\,000~\mathrm{kg/cm^2}.
\end{equation}

\subsection{Parámetros geotécnicos}
Los parámetros del suelo de fundación corresponden a los asignados para el sector en estudio y se resumen en la Tabla~\ref{tab:geotecnicos}. Estos valores corresponden a tensiones admisibles a nivel de sello de fundación, sin que corresponda descontarles el peso del suelo dispuesto sobre dicho nivel.

\begin{table}[H]
\centering
\caption{Parámetros geotécnicos de diseño.}
\label{tab:geotecnicos}
\begin{tabular}{lcc}
\hline
Parámetro & Símbolo & Valor \\
\hline
Presión admisible estática & $\sigma_{adm,e}$ & 70~tonf/m$^2$ \\
Presión admisible sísmica & $\sigma_{adm,s}$ & 93~tonf/m$^2$ \\
Coeficiente de roce suelo-fundación & $\mu$ & 0,40 \\
Factor de seguridad mínimo al deslizamiento & $FS_d$ & 1,50 \\
Profundidad de sello de fundación & $D_f$ & 2,50~m \\
Peso específico del suelo & $\gamma_s$ & 1,836~tonf/m$^3$ \\
Peso específico del hormigón & $\gamma_h$ & 2,50~tonf/m$^3$ \\
\hline
\end{tabular}
\end{table}

\subsection{Geometría de las zapatas y criterios de diseño}
Se adopta una altura uniforme $H = 1,20$~m para todas las zapatas, con un recubrimiento de 0,05~m, lo que entrega una altura útil $d = 1,15$~m. Las verificaciones estructurales se realizan conforme a ACI 318-19, empleando factores de reducción de resistencia $\phi = 0,90$ para flexión y $\phi = 0,75$ para corte. Para el diseño estructural de la zapata se emplea un factor de mayoración de cargas de 1,4.

La armadura mínima por retracción y temperatura se determina con una cuantía de 0,0018 sobre la sección bruta:
\begin{equation}
A_{s,\min} = 0,0018 \cdot b \cdot H = 0,0018 \cdot 100~\mathrm{cm} \cdot 120~\mathrm{cm} = 21,6~\mathrm{cm^2/m}.
\end{equation}

Para los muros y columnas con zapata centrada, se establece como criterio que la zona comprimida bajo la fundación abarque al menos el 80\% de su longitud, limitando la magnitud del despegue admisible entre la zapata y el suelo. Este criterio no resulta aplicable a los muros perimetrales con zapata excéntrica, dado que la fracción de base comprimida en ese caso queda determinada por la condición geométrica de borde y no por la magnitud de la solicitación; su tratamiento se detalla en la sección 5.3.